% ======================================================================
% Chapter 1 -- Part B: the four building blocks; decomposing the quote desk
% ======================================================================

\section{Four Building Blocks}
\label{sec:blocks}

The book's subtitle names three technologies: RAG, agents and MCP. In practice, every solution in this book also uses a fourth building block that is easy to forget because it is not new: plain code. We take the four in order, from the most predictable to the least.

\subsection{Plain code and rules}
\label{sec:block-code}

Plain code is any logic a programmer writes and a test can check exactly: a policy threshold, a currency conversion, a database query, a schema check, a routing rule. It is fast, nearly free to run, perfectly repeatable and easy to audit. It is the right block whenever the rule can be written down and the inputs are structured.

The common mistake in agentic designs is not using too little AI; it is putting rules inside prompts. When Team A wrote ``never approve more than 15\% without manager sign-off'' in the system prompt, they turned a rule that code enforces with certainty into a request that a model follows most of the time. A rule in a prompt cannot be unit-tested, cannot be audited line by line, and can be argued around by a cleverly worded input. A rule in code can be all three.

\begin{designnote}
If a requirement contains the words \emph{never}, \emph{always}, \emph{must} or a number, start by asking whether it belongs in code. Prompts can \emph{explain} a rule to the model so it behaves sensibly; code should \emph{enforce} it.
\end{designnote}

\subsection{Retrieval-augmented generation (RAG)}
\label{sec:block-rag}

A language model knows what was in its training data, which is general, dated and silent about your company. \keyterm{Retrieval-augmented generation} supplies the missing knowledge at the moment of use: the system searches a store of documents or records for passages relevant to the request and places them in the model's context before the model answers \citep{lewis2020rag}. At Harbor, the discount policy, the product catalog notes and each customer's contract are knowledge that must come from retrieval, not from the model's memory.

RAG is the right block when the system needs \emph{knowledge it cannot be trained on}: private, frequently changing, or too large to put in every prompt. It is often the wrong block for \emph{structured facts with an exact answer}. The current list price of part \texttt{BV-316-3} is a database lookup, not a search problem; using semantic search for it invites the system to return the price of a similar part.

Designing a RAG component means answering a short list of questions, each of which comes back in later chapters:
\begin{itemize}[itemsep=1pt]
  \item \emph{Sources:} which documents, owned by whom, in what formats?
  \item \emph{Freshness:} how quickly must a change in a source reach the answers? Minutes, a day, a quarter?
  \item \emph{Permissions:} may every user see every document? (At Harbor, a sales representative may see the policy but not another region's contract terms.)
  \item \emph{Granularity:} how are documents split into passages, and does a split ever cut a rule in half?
  \item \emph{Grounding:} must every answer cite its source and version, and what happens when nothing relevant is found?
  \item \emph{Quality:} how will you measure whether the right passages were retrieved, separately from whether the final answer was good?
\end{itemize}
Chapter~\ref{ch:case-policy}, the first case, goes deep on these questions.

\subsection{The agent}
\label{sec:block-agent}

The third block is the agent as defined in Section~\ref{sec:autonomy}: a model that chooses its next action in a loop. The pattern most frameworks implement is some form of \emph{reason, act, observe}: the model reasons about the goal, calls a tool, reads the result, and repeats until it decides to stop \citep{yao2023react}. In ADK the loop is built in; you supply the model, the instructions, and the tools.

An agent is the right block when a step needs \emph{judgment over a path that cannot be fixed in advance}: understanding an unstructured request, deciding which of several lookups is needed, reconciling conflicting information, or writing text for a person to read. At Harbor, reading a rambling email from a sales representative and turning it into a structured quote request needs judgment. So does writing the two-sentence justification that a manager reads before approving a 17\% discount.

Agents fail in characteristic ways, and a designer should expect each one:
\begin{itemize}[itemsep=1pt]
  \item \emph{Wrong tool or wrong arguments:} the agent looks up the wrong customer because two have similar names.
  \item \emph{Unsupported claims:} the agent states a contract term that is not in any retrieved document.
  \item \emph{Loops and runaway cost:} the agent keeps calling tools without converging.
  \item \emph{Instruction drift:} a long conversation or an adversarial input pushes the agent away from its instructions.
  \item \emph{Overreach:} the agent takes an action, such as sending an email, that should have waited for a person.
\end{itemize}
Each failure has a design response: fewer and better-shaped tools, required citations, step and cost limits, guardrails, and human checkpoints. The cases teach them one at a time.

\subsection{The Model Context Protocol (MCP)}
\label{sec:block-mcp}

An agent is only as useful as the actions it can take and the context it can reach. Before MCP, every framework had its own way of describing a tool, and every integration with a company system was written again for each framework. The \keyterm{Model Context Protocol} is an open standard that separates the two sides \citep{mcp2026spec}. A \emph{server} exposes capabilities; a \emph{client}, running inside a \emph{host} application such as an ADK agent, a desktop assistant or an IDE, discovers and uses them. A server written once can serve any MCP-compatible client.

An MCP server can offer three kinds of capability, and choosing among them is itself a design decision:
\begin{description}[leftmargin=1em,itemsep=2pt]
  \item[Tools] are actions the model may decide to invoke, with typed inputs and outputs. Examples are a tool that returns the contract price for a customer and a SKU, or one that creates a quote draft. A tool may read or change the world.
  \item[Resources] are read-only context that the host application can attach, addressed by URI: the current discount policy, a product data sheet. Resources are chosen by the application or the user rather than by the model.
  \item[Prompts] are reusable, parameterized message templates that the server offers to users or applications: ``draft a quote justification for manager review.''
\end{description}
The usual transports are \emph{stdio}, for a server running as a local process beside the client, and \emph{Streamable HTTP}, for a remote server. The protocol is versioned by date. The revision current as of this writing, dated 2026-07-28, made the core protocol stateless: each request carries what it needs, so a remote MCP server can run behind an ordinary load balancer like any other web service \citep{mcp2026spec}. For a designer, this means an MCP server can be deployed, scaled and secured with the same tools as the rest of the company's services.

Why does a protocol belong in a book about design? Because MCP gives us a natural \emph{boundary}. Deciding what goes into an MCP server, and what each tool's contract looks like, is the same kind of decision as deciding the interface of a microservice. A good MCP server is owned by the team that owns the underlying system, exposes a small number of task-shaped tools rather than a thin wrapper over every API endpoint, enforces its own permissions, and can be tested without any model in the loop. A bad one exposes \texttt{run\_sql(query)} and hopes for the best. Chapter~\ref{ch:case-support} is devoted to tool design.

A related protocol, \keyterm{Agent2Agent (A2A)}, addresses a different boundary: communication between agents that are built and run by different teams or vendors. MCP connects an agent to tools and context; A2A connects an agent to other agents. We use A2A in the capstone, Chapter~\ref{ch:capstone}.

\subsection{Choosing among the blocks}

Table~\ref{tab:blocks} summarizes the four blocks. Keep it beside you when you decompose a process; the right-hand column is often the fastest way to catch a poor assignment.

\begin{table}[htbp]
\centering
\small
\caption{The four building blocks of an agentic solution.}
\label{tab:blocks}
\begin{tabularx}{\textwidth}{@{}>{\raggedright\arraybackslash}p{2.1cm}>{\raggedright\arraybackslash}X>{\raggedright\arraybackslash}X>{\raggedright\arraybackslash}X@{}}
\toprule
Block & Supplies & Use it when & Warning sign you chose wrong \\
\midrule
Plain code & Certainty: rules, calculations, routing, validation & The rule can be written down and inputs are structured & You are writing ``never'' or a threshold into a prompt \\
RAG & Knowledge: private, changing or large text & Answers depend on documents the model was not trained on & You are using semantic search to find an exact record \\
Agent & Judgment: interpreting, choosing, writing & The path cannot be fixed in advance or inputs are unstructured & The agent follows the same three steps every time \\
MCP server & Action and context: access to systems of record & The system must read from or write to company systems & A tool mirrors a raw API, or one tool can do anything \\
\bottomrule
\end{tabularx}
\end{table}

Two further components appear in every design but are not building blocks in the same sense. A \keyterm{human checkpoint} is a point where the system stops and a person approves, edits or rejects its work. A \keyterm{guardrail} is a check, usually in code, that inspects inputs or outputs and blocks or repairs what violates a policy. Both are first-class elements of the notation in Chapter~\ref{ch:toolkit}.

\subsection{How the blocks work together}
\label{sec:blocks-together}

Table~\ref{tab:blocks} treats the blocks one at a time, but a real solution uses them together. Let's look at an example. Figure~\ref{fig:integrated} shows a quote assistant for Harbor's pricing desk. A sales representative sends a request in plain words: 40 valves for Delta Water at 18\% off. The \texttt{quote\_desk} agent reads the request and decides, step by step, which tool to call next; that judgment is the only part a model performs. Whether 18\% is allowed is not left to the model: a plain-code function, \texttt{check\_discount}, enforces the 15\% cap. Everything the agent needs from outside travels through one MCP client. Retrieval over the policy and contract text supplies the clause that applies and its version. The ERP server returns list price, unit cost, stock and contract terms. The CRM server records the finished quote as a draft. The last step belongs to a person: the representative reviews and sends the quote, after a manager has approved any discount above the cap.

Notice how each block does only the job Table~\ref{tab:blocks} assigns to it. The model reads and decides, code enforces, retrieval explains, and MCP servers connect to the systems of record. Section~\ref{sec:decompose} shows how to arrive at this assignment from the business process, one step at a time, and Section~\ref{sec:adk} builds the agent and the pricing server.

% Figure: the four building blocks working together in Harbor's quote assistant.
% Uses the color names of fig:quote-sketch (agentfill, codefill, mcpfill, ragfill) from
% ../common/preamble. TikZ libraries: positioning, shapes.geometric, arrows.meta, fit, backgrounds.
\begin{figure}[htbp]
\centering
\begin{tikzpicture}[
  node distance=7mm and 8mm,
  box/.style={draw, rounded corners=2pt, align=center, font=\footnotesize, minimum height=10mm},
  agent/.style={box, fill=agentfill, text width=4.6cm},
  code/.style={box, fill=codefill, text width=2.2cm},
  mcp/.style={box, fill=mcpfill, text width=2.8cm},
  rag/.style={box, fill=ragfill, shape=cylinder, shape border rotate=90, aspect=0.18,
              rounded corners=0pt, text width=2.5cm},
  io/.style={box, fill=white, text width=2.1cm},
  human/.style={box, fill=humanfill, shape=rectangle, double, text width=2.1cm},
  sor/.style={box, fill=white, dashed, text width=2.6cm},
  arr/.style={-{Stealth[length=2mm]}, thick},
  biarr/.style={{Stealth[length=2mm]}-{Stealth[length=2mm]}, thick},
  lab/.style={font=\scriptsize, fill=white, inner sep=1pt}
]
% inside the agent
\node[agent] (llm) {LLM loop (agent)\\{\scriptsize reads the request, picks a tool, reads the result, repeats}};
\node[code, below left=8mm and -2.2cm of llm] (disc) {{\ttfamily\scriptsize check\_discount}\\{\scriptsize 15\% cap (code)}};
\node[box, fill=mcpfill, text width=2.2cm, below right=8mm and -2.2cm of llm] (client) {MCP client};

% agent boundary
\begin{scope}[on background layer]
  \node[draw, dashed, rounded corners=4pt, inner sep=3mm, fit=(llm)(disc)(client),
        label={[font=\footnotesize\bfseries]above:{\ttfamily quote\_desk} agent (ADK)}] (agentbox) {};
\end{scope}

% request and result
\node[io, left=9mm of llm] (rep) {Sales rep\\{\scriptsize ``40 valves for Delta Water, 18\% off''}};
\node[human, right=9mm of llm] (send) {Rep reviews and sends\\{\scriptsize manager first if $>$15\%}};

% the MCP bus and servers
\coordinate (bus) at ([yshift=-7mm]agentbox.south -| client);
\node[mcp, anchor=north] (erp) at ([yshift=-15mm]agentbox.south) {ERP MCP server\\{\scriptsize\ttfamily get\_item\_pricing}\\{\scriptsize\ttfamily get\_contract\_terms}};
\node[rag, anchor=north east] (kb) at ([xshift=-7mm]erp.north west) {Policy \&\\contract text (RAG)\\{\scriptsize cites policy version}};
\node[mcp, anchor=north west] (crm) at ([xshift=7mm]erp.north east) {CRM MCP server\\{\scriptsize\ttfamily create\_quote\_draft}};
\node[sor, below=7mm of erp] (erpdb) {ERP\\(system of record)};
\node[sor, below=7mm of crm] (crmdb) {CRM\\(system of record)};

% arrows
\draw[arr] (rep) -- (llm);
\draw[arr] (llm) -- (send);
\draw[biarr] (llm.south -| disc) -- (disc);
\draw[biarr] (llm.south -| client) -- (client);
\draw[thick] (client.south) -- (bus);
\draw[thick] (bus -| kb) -- (bus -| crm);
\draw[biarr] (bus -| kb) -- (kb.north);
\draw[biarr] (bus -| erp) -- (erp.north);
\draw[biarr] (bus -| crm) -- (crm.north);
\node[lab, above right=0.5mm and 1mm of bus] {MCP};
\draw[arr, dashed] (erp) -- (erpdb);
\draw[arr, dashed] (crm) -- (crmdb);

% legend: the four building blocks
\begin{scope}[
  swatch/.style={draw, rounded corners=1pt, minimum width=5mm, minimum height=3.5mm, inner sep=0pt},
  key/.style={font=\footnotesize, anchor=west, inner sep=1pt}
]
  \coordinate (leg) at ([yshift=-8mm]current bounding box.south -| erp);
  \node[swatch, fill=codefill] (s1) at ([xshift=-41mm]leg) {};
  \node[key, right=1.5mm of s1] (k1) {Plain code};
  \node[swatch, fill=agentfill, right=6mm of k1] (s2) {};
  \node[key, right=1.5mm of s2] (k2) {Agent};
  \node[minimum width=4.5mm, minimum height=6mm, inner sep=0pt, right=6mm of k2] (s3) {};
  \filldraw[fill=ragfill] (s3.north west) ++(0,-0.8mm) -- ++(0,-4.4mm)
    arc[start angle=180, end angle=360, x radius=2.25mm, y radius=0.8mm] -- ++(0,4.4mm)
    arc[start angle=0, end angle=180, x radius=2.25mm, y radius=0.8mm] -- cycle;
  \draw ([yshift=-0.8mm]s3.north west)
    arc[start angle=180, end angle=360, x radius=2.25mm, y radius=0.8mm];
  \node[key, right=1.5mm of s3] (k3) {RAG};
  \node[swatch, fill=mcpfill, right=6mm of k3] (s4) {};
  \node[key, right=1.5mm of s4] (k4) {MCP};
  \begin{scope}[on background layer]
    \node[draw, gray, rounded corners=2pt, inner sep=2mm, fit=(s1)(k4)] {};
  \end{scope}
\end{scope}
\end{tikzpicture}
\caption{The four building blocks working together in Harbor's quote assistant. The agent decides which step comes next; plain code enforces the discount cap; RAG supplies the policy text and its version; MCP servers give access to the ERP and the CRM. A person approves before anything leaves the company.}
\label{fig:integrated}
\end{figure}


% ----------------------------------------------------------------------
\section{Decomposing a Process: The Pricing Desk}
\label{sec:decompose}

We now apply the blocks to Harbor's pricing desk. The method is the first half of the design method that Chapter~\ref{ch:toolkit} formalizes: write down the current process step by step, then decide for each step whether it needs judgment, knowledge, action, a fixed rule---or a person.

Table~\ref{tab:quote-decomp} shows the result. The current process is in the first two columns, taken from interviews with the pricing analysts. The last two columns are design decisions.

\begin{table}[htbp]
\centering
\small
\caption{Decomposing Harbor's quote process. The ``needs'' column drives the choice of block.}
\label{tab:quote-decomp}
\begin{tabularx}{\textwidth}{@{}c>{\raggedright\arraybackslash}X>{\raggedright\arraybackslash}p{2.2cm}>{\raggedright\arraybackslash}X@{}}
\toprule
\# & Step today & Needs & Block in the design \\
\midrule
1 & Read the rep's email or chat message; identify customer, items, quantities, requested discount & Judgment & Agent (Level 1 extraction into a fixed schema) \\
2 & Check the request is complete; ask the rep for anything missing & Rule & Code: schema validation; agent writes the follow-up question \\
3 & Look up list price, unit cost and stock for each item & Action & MCP server over the ERP: \texttt{get\_item\_pricing} \\
4 & Look up the customer's contract price and terms & Action + knowledge & MCP tool for structured contract fields; RAG over contract text for special clauses \\
5 & Check the requested discount against policy & Knowledge + rule & Code enforces cap and margin floor; RAG cites the policy section and version \\
6 & If above 15\% or below margin floor, get manager approval & Person & Human checkpoint with a model-written justification \\
7 & Write the quote and the cover note & Judgment & Agent drafts; code fills all numbers from steps 3--5 \\
8 & Record the quote in the CRM and send it & Action & MCP server over the CRM; sending waits for the rep's click \\
\bottomrule
\end{tabularx}
\end{table}

Three decisions in the table are worth pausing on, because they are the decisions Team A did not make.

\begin{itemize}[itemsep=2pt]
  \item \emph{The model never computes a price.} In step 7 the agent writes the words, but every number in the quote is filled in by code from the results of steps 3 to 5. This removes an entire class of error---an arithmetic slip or a hallucinated figure on a legally binding document---at almost no cost.
  \item \emph{Policy is enforced in code and explained by retrieval.} The cap and the margin floor are numbers in a configuration file, tested and versioned. The policy document is still retrieved, so the manager sees which section and version applied, and so the system notices when the document and the configuration disagree.
  \item \emph{Two people stay in the loop for different reasons.} The manager approves exceptions because the business wants control over margin. The sales representative sends the quote because the customer relationship is theirs. Neither checkpoint is there because the model is untrustworthy in general; each is there because a specific decision belongs to a specific person.
\end{itemize}

Figure~\ref{fig:quote-sketch}, Team B's diagram from Section~\ref{sec:two-demos}, is this table drawn as a first sketch. It is deliberately informal; Chapter~\ref{ch:toolkit} replaces it with the book's standard notation, and Chapter~\ref{ch:business} attaches KPIs to each node. Even informal, the sketch makes the design arguable: a reviewer can point at any box and ask why it is there and what happens when it fails.

\subsection{What the sketch does not yet say}

A sketch is the beginning of a design, not the end. Before Team B could defend theirs, they had to answer questions the sketch leaves open---the same questions every case in this book asks:
\begin{itemize}[itemsep=1pt]
  \item \emph{Alternatives.} Would a single agent with all the tools (Team A's design) be cheaper to build and good enough? Would a pure workflow with no agent at all handle 90\% of requests?
  \item \emph{Business value.} If turnaround falls from 26 hours to 4, how many more deals does Harbor win, and what is that worth? What does each quote cost to produce, in model calls and in analyst minutes?
  \item \emph{Cost of errors.} A quote below the margin floor loses money on every unit; a quote above the customer's contract price damages trust. Which error is worse, and does the design guard hardest against the worse one?
  \item \emph{Production.} How is the system evaluated before launch, monitored after, and rolled back if something goes wrong?
\end{itemize}

A first pass at unit economics already shows why these questions matter. Let $m$ be the analyst minutes per quote, $w$ the loaded cost per analyst minute, $c$ the model and infrastructure cost per quote, and $r$ the share of quotes that still need a manager's review, taking $m_r$ minutes each. The cost of producing one quote is
\begin{equation}
  C_{\text{quote}} = c + w \bigl( m + r\, m_r \bigr).
  \label{eq:quote-cost}
\end{equation}
Today, with no AI, $c = 0$ and $m$ is about 35 minutes. Team B's design aims to cut $m$ to about 5 minutes (a pricing analyst spot-checks the draft) while adding a model cost $c$ of a few cents. Even if $c$ were ten times higher than planned, it would remain small next to the analyst time saved. The real risk to the business case is elsewhere: in $r$, if the system routes too many quotes to managers, and in the cost of errors, which Equation~\eqref{eq:quote-cost} leaves out entirely. Chapter~\ref{ch:business} adds that term and shows how to estimate it.
